\documentclass[10pt,a4paper]{article}
\usepackage{fontspec,amsmath,amssymb,graphicx}
\usepackage[margin=18mm]{geometry}
\setmainfont{Times New Roman}
\setlength{\parindent}{0pt}
\setlength{\parskip}{5pt}
\allowdisplaybreaks
\sloppy
\emergencystretch=2em
\begin{document}
\section*{BT\_5\_4\_3\_13 — Lời giải bổ sung}
Ký hiệu $B$ trong đề dùng $|x|\le1$: nếu $u$ liên tục trên quả cầu đóng thì Heine--Cantor buộc liên tục đều, nên yêu cầu như vậy không thể thỏa. Với ý định là quả cầu mở, trong đĩa lấy
\[u(x,y)=\operatorname{Re}\exp\left[-\frac{1+x+iy}{1-x-iy}\right].\]
Hàm giải tích vì cực $z=1$ nằm ngoài miền; $\operatorname{Re}(1+z)/(1-z)>0$ nên $|u|\le1$. Khi $z\to1$ theo bán kính, $u\to0$; theo các đường có phần thực thương tiến về một hằng và phần ảo thay đổi, $u$ có nhiều giá trị giới hạn. Do đó không liên tục đều. Có thể mở rộng ví dụ sang $n\ge2$ bằng cách không phụ thuộc các tọa độ còn lại.\end{document}